% vim: set tw=78 tabstop=4 shiftwidth=4 aw ai:
\documentclass{beamer}

\usepackage[utf8x]{inputenc}		% diacritice
\usepackage[english]{babel}
\usepackage{color}			% highlight
\usepackage{alltt}			% highlight

% highlight; comment this out in case you don't input code source files
\usepackage{code/highlight}		% highlight

\usepackage{hyperref}			% folosiți \url{http://...}
					% sau \href{http://...}{Nume Link}
\usepackage{verbatim}

\mode<presentation>
{ \usetheme{Berlin} }

% Încărcăm simbolurilor Unicode românești în titlu și primele pagini
\PreloadUnicodePage{200}

% Arătăm numărul frame-ului
\newcommand{\frameofframes}{/}
\newcommand{\setframeofframes}[1]{\renewcommand{\frameofframes}{#1}}

\setframeofframes{of}
\makeatletter
\setbeamertemplate{footline}
  {%
    \begin{beamercolorbox}[colsep=1.5pt]{upper separation line foot}
    \end{beamercolorbox}
    \begin{beamercolorbox}[ht=2.5ex,dp=1.125ex,%
      leftskip=.3cm,rightskip=.3cm plus1fil]{author in head/foot}%
      \leavevmode{\usebeamerfont{author in head/foot}\insertshortauthor}%
      \hfill%
      {\usebeamerfont{institute in head/foot}\usebeamercolor[fg]{institute in head/foot}\insertshortinstitute}%
    \end{beamercolorbox}%
    \begin{beamercolorbox}[ht=2.5ex,dp=1.125ex,%
      leftskip=.3cm,rightskip=.3cm plus1fil]{title in head/foot}%
      {\usebeamerfont{title in head/foot}\insertshorttitle}%
      \hfill%
      {\usebeamerfont{frame number}\usebeamercolor[fg]{frame number}\insertframenumber~\frameofframes~\inserttotalframenumber}
    \end{beamercolorbox}%
    \begin{beamercolorbox}[colsep=1.5pt]{lower separation line foot}
    \end{beamercolorbox}
  }
\makeatother

\setbeamertemplate{navigation symbols}{}%remove navigation symbols

\title[My Super Bachelor Thesis]{My Super Bachelor Thesis}
\subtitle{Bachelor Thesis Session -- July 2010}
\institute{Faculty of Automatic Control and Computers,\\
	University POLITEHNICA of Bucharest}
\author[Ion Ionescu]{Ion Ionescu\\
	Supervisor: Andrei Antonescu}
\date{July 9, 2010}

\begin{document}

% Slide-urile cu mai multe părți sunt marcate cu textul (cont.)
\setbeamertemplate{frametitle continuation}[from second]

\frame{\titlepage}

\frame{\tableofcontents}

% NB: Secțiunile nu sunt marcate vizual, ci doar apar în cuprins
\section{Context}

% Titlul unui frame se specifică fie în acolade, imediat după \begin{frame},
% fie folosind \frametitle
\begin{frame}{TODO}
	\begin{itemize}		% Just like normal LaTeX
		\item TODO
	\end{itemize}
\end{frame}

\section{Architecture and Design}

\begin{frame}{TODO}
	\begin{itemize}
		\item Basic formula
			\begin{beamerboxesrounded}[lower=block body,shadow=true]{}
				\texttt{\#define MAX(a, b)   ((a) > (b) ? (a) : (b))}
			\end{beamerboxesrounded}
		\item TODO
	\end{itemize}
\end{frame}

\section{Implementation}

\begin{frame}{TODO}
	\begin{itemize}
		\item TODO
		\item TODO
	\end{itemize}
\end{frame}

\begin{frame}{Code snippets}
	\documentclass{propsoa.cs.pub.ro}

\title{Project Name}
\subtitle{Project Proposal, SOA/OS -- Autumn 2010}
\version{1.1/\today}
\indexterms{project scope}
\keywords{project/paper specific topics}
\teamsize{number of students in team}

\begin{document}
\maketitle

\section{Project Description}

The scientific context in which the project is developed must be presented.
The purpose and objective of the project must be properly stated.

\section{Objectives}

This project aims to implement a software application that will meet the
following requirements:
\begin{itemize}
	\item{Objective 1}
	\item{Objective 2}
	\item{Objective 3}
\end{itemize}

\section{Bibliography}

list of papers, books, links

\section{Prerequisites}

Fields of study, technologies.

\section{Other}

(other useful information, etc. --- section may be missing)

\end{document}

\end{frame}
	
\section{Results}

\begin{frame}{TODO}
	\begin{itemize}
		\item TODO
		\item TODO
	\end{itemize}
\end{frame}

\section{Questions}

\begin{frame}{Questions}
  \begin{columns}
    \begin{column}[t]{0.5\textwidth}
      \begin{itemize}
        \item keyword1
        \item keyword2
        \item keyword3
        \item keyword4
        \item keyword5
      \end{itemize}
    \end{column}
    \begin{column}[c]{0.5\textwidth}
      \begin{figure}
        \includegraphics[scale=0.3]{img/question-mark}
      \end{figure}
    \end{column}
  \end{columns}
\end{frame}

\end{document}
